\section{Introduction}

Coding agents are increasingly used to automate software-engineering tasks
through extended interactions with development
environments~\cite{sweagent,openhands,autocoderover,hyperagent}. In these
workflows, agents edit code and execute commands, using the resulting
observations to decide how to proceed~\cite{react,toolformer}. Yet an
incorrect edit or command can disrupt the task and cause subsequent actions
to build on an erroneous result. Worse still, the effects of the mistake can
persist even after it is recognized, leaving the agent to repair affected work
before making further progress~\cite{agentrewind,planarian}. Rollback is
therefore essential for removing these effects and restoring a valid state
from which the agent can continue~\cite{try-osdi26,branchfs,crab}.

To support rollback, existing systems stage effects or restore checkpoints.
In particular, effect-staging systems isolate filesystem changes for later
application or discard~\cite{try-osdi26,branchfs,cordon}. When selecting
file effects to discard, the caller also has to identify dependent calls
and observations to remove. By comparison, checkpoint/restore systems save state for later
restoration~\cite{rollbacksurvey,dmtcp,waypoint,deltabox}. Building on
these mechanisms, recent agent runtimes select checkpoint state from turn effects~\cite{crab}, restore
consistent environment versions~\cite{planarian}, or recover aligned
filesystem and agent-context checkpoints while retaining lessons~\cite{agentrewind}.
In the 2025 Stack Overflow survey, 66\% of respondents ($n=31{,}476$)
reported ``almost right, but not quite'' AI solutions~\cite{stackoverflow-survey-2025};
45.1\% of merged requests in a study of 567 Claude Code-assisted pull
requests underwent revision~\cite{watanabe-agentic-prs}.
Users may want to roll back selected files while retaining other work,
as reflected in file-level keep/undo controls~\cite{vscode-review-edits}.
Checkpoint and branch rollback can discard that work; file-level undo
alone does not remove dependent effects or agent observations.

Our motivational experiments (Section~\ref{sec:motivation}) show that
the DeltaBox recovery scope discards 56.90\% of valid work on average across 400
fixed trajectories, even with exact fault information. Moreover, dependent
calls interleave with independent work in 86.8\% of 946 complete clean agent
traces. These observations show that existing recovery policies discard
substantial valid work. Our goal is to retain independent progress while
undoing affected file changes and removing the associated tool results
from agent context.
To this end, we present \system, a causal recovery runtime for coding
agents. The key idea is to \textbf{let dependencies among tool calls
determine the recovery boundary across filesystem state and agent context}.
Given the files to roll back, \system traces file dependencies and
tool-result reuse from their changes to derive one rollback set.
An isolated trial checks the expected file state and command/result residue
before live recovery. The runtime restores affected file objects and
removes the corresponding calls and results from agent context while preserving
independent progress. Consequently, recovery can preserve independent calls
within and across turns and combine earlier file versions with later valid
updates, even when that combination never existed at a checkpoint.

However, designing a filesystem that supports causal rollback poses three
key challenges.
1) Errors can propagate through both file accesses and reused tool results.
\system combines these dependencies to select affected calls and checks
the resulting rollback state before resuming execution.
2) The desired file state may be absent from every checkpoint. To
reconstruct it, \system restores selected object versions, metadata, and
deletion state in OverlayFS while retaining independent updates.
3) A single agent turn can contain both affected and independent calls.
Accordingly, \system uses shared call identifiers to remove affected
calls and results from the next model request while retaining independent
pairs within the same turn.

We implement \system on Linux OverlayFS and integrate it with SWE-agent
and Codex. In controlled
experiments on 400 workloads, with three
repetitions per policy, \system removes all affected calls while retaining
all independent work. Beyond these controlled workloads, \system reduces token consumption
on real-world SWE-Bench Lite~\cite{swebench} and Terminal-Bench~\cite{tbench}
tasks relative to DeltaBox by 17.55--23.45\% and 25.39--44.69\%,
respectively, on matched successful task pairs across the two agents.
On the same real tasks, the breakdown shows lower agent-execution and
setup/replay time relative to DeltaBox.

This paper makes three contributions:
\begin{itemize}
  \item We quantify the valid work lost when recovery boundaries include
        calls independent of an error, and show that dependent calls
        commonly interleave with independent work in real agent traces
        (Section~\ref{sec:motivation}).
  \item We design and implement \system, which combines file dependencies
        and tool-result reuse to select and validate rollback sets. The
        selected set guides selective object restoration and updates
        agent context, preserving independent calls within and across turns
        (Section~\ref{sec:design}).
  \item We evaluate recovery correctness, preserved work, runtime overhead,
        and model cost on controlled and real-world workloads
        (Section~\ref{sec:eval}).
\end{itemize}
